\scanpage{191}
\mspage{743}

structure -- linked to a class of discrete subgroups commensurable with one another -- of \deleted{meta-$L$-}\inserted{$L$-$T$-}groups. (where $L$ is the initial of ``loops'', and $T$ of ``Teichmüller'') -- these are sorts of \unsure{mixtures} of Teichmüller groups and $L$-groups. To illustrate this point, let us start from an extension
\diagram{(6)\qquad
\begin{tikzcd}[ampersand replacement=\&, column sep=large]
1 \arrow[r] \& \Pi \arrow[r] \& \mathcal{G} \arrow[r] \& \mathcal{T} \arrow[r] \& 1
\end{tikzcd}
}
where $\Pi$ is equipped with an $L$-structure invariant under $\mathcal{G}$, then $\mathcal{G}$ is an $L$-$T$-group admitting (if $\Pi$ is itself a \deleted{\unsure{$L$-}}group with loops) $\Pi$ as a subgroup of definition. Starting from the group with loops $\Pi$, one can take for example
\[
(7)\qquad
\left\{
\begin{array}{l}
\mathcal{G}=\mathrm{Aut}_{L}(\Pi)\\
\mathcal{T}=\mathrm{Autext}_{L}(\Pi)
\end{array}
\right.
\qquad\text{whence}
\,,
\]

\scanpage{192}

and thus the \unsure{ambient} group $\mathcal{G}$, an extension of the Teichmüller group $\mathcal{T}$ by the group with loops $\Pi$, appears as an $LT$-group.

\deleted{\unsure{2}}

\scanpage{193}
\mspage{744}

2$^\circ$) \uline{Profinite case}

The extension of the notion of an $L$-group or of an $LT$-group \inserted{$\Pi$}, in terms of \uline{germs} of outer homomorphisms
\diagram{(8)\qquad
\begin{tikzcd}[ampersand replacement=\&]
k_i^{\unsure{\natural|\sharp}}\colon T \arrow[r] \& \Pi
\end{tikzcd}\,,
}
is essentially evident. Here is the most important example for us. \deleted{\unsure{Let}}

Let first $k$ be a field \deleted{algebraically closed of characteristic $0$}, and let $M$ be an \deleted{\unsure{0-connected}} \deleted{\unsure{0-connected}} algebraic multiplicity over $k$, which is \inserted{0-connected,} of finite type over $k$, \deleted{smooth} separated$/k$ and smooth \unsure{in|of} relative dimension $1$. One should be able to \unsure{prove} then that $M$ can be described by a smooth and quasi-projective algebraic curve $U$ over $k$, a finite group \deleted{$\unsure{K}$} $G$ which operates on $U$, \deleted{and} as
\[
(9)\qquad M=(U,G)=
\begin{array}{l}
\inserted{\text{(small)}}\text{ Ringed topos of}\\
\text{étale sheaves on}\\
U\text{ where }G\text{ operates}\ldots
\end{array}
\]

\scanpage{194}

\deleted{\unsure{or}} -- More intrinsically, it should be true that there exists a finite étale covering of $M$ (which one can suppose Galois, with group $G$) \deleted{\unsure{and|such}} which is \unsure{an|a} \deleted{algebraic curve}\inserted{scheme} -- which will necessarily be \unsure{quasi-projective|quasi-projective}, smooth, of relative dimension $1$ over $k$. Let us suppose in any case that this is so, choose a geometric base point $a$ of $M$, and consider
\[
(10)\qquad \Pi=\Pi_1(M,a)\simeq\Pi_1(\unsure{U|M},G;a)\,.
\]

We shall henceforth suppose $k$ of characteristic zero. Let us first suppose $k$ algebraically closed. Using the exact sequence
\diagram{(11)\qquad
\begin{tikzcd}[ampersand replacement=\&, column sep=large]
1 \arrow[r] \& \underbrace{\Pi_1(U,a)}_{\Pi_0} \arrow[r] \& \underbrace{\Pi_1(\unsure{U|M},G;a)}_{\Pi} \arrow[r] \& G \arrow[r] \& 1
\end{tikzcd}
}
and the natural \deleted{\unsure{$L$}} loop structure on $\Pi_0$, one finds an $L$-structure on $\Pi=\Pi_1(\unsure{M},a)$, which manifestly does not

\scanpage{195}
\mspage{745}

depend on the choice of the \inserted{finite} Galois étale covering $U$ of $M$, and on the chosen lifting of the geometric point $a$ of $M$ to a geometric point of $U$. We shall always \inserted{tacitly} equip $\Pi=\Pi_1(M,a)$ with this $L$-structure.

If again $k$ is an arbitrary field (the most interesting case being the one where $k$ is an extension of finite type of $\mathbb{Q}$), we shall have an exact sequence
\diagram{(12)\qquad
\begin{tikzcd}[ampersand replacement=\&, column sep=large]
1 \arrow[r] \& \Pi_1(M_{\overline{k}},a) \arrow[r] \& \Pi_1(M,a) \arrow[r] \& \Gamma_{\overline{k}/k} \arrow[r] \& 1
\end{tikzcd}
}
and $\Pi_1(M_{\overline{k}},\unsure{a|o})$ appears as equipped with an $L$-structure invariant under $\Pi_1(M,\unsure{a|o})=\Pi$, so that $\Pi$ appears as a $TL$-group.

I am particularly interested in the ``universal'' case of type $(g,\nu)$ -- but the \inserted{natural} base is not then the

\scanpage{196}
\mspage{\deleted{\unsure{745}}}

spectrum of a field $k$, but \unsure{the|the} modular \deleted{\unsure{scheme}} \inserted{multiplicity} $M_{g,\nu}$. [To tell the truth, for the discrete $L$-group structure $\Pi_0$, playing the role of reference group, one \deleted{can} \inserted{should be able to} introduce a corresponding modular \deleted{\unsure{scheme}} \inserted{multiplicity} over $\mathbb{Q}$, and even over $\mathbb{Z}$. The fact which risks causing \unsure{conceptual} difficulties for such a construction, it is that a \inserted{``modular''} multiplicity (over $\mathbb{Q}$, or over $\mathbb{Z}$\ldots) describes in principle a fibered \uline{category} --- such as that of relative curve schemes of type $(g,\nu)$ --- whereas the relative $\unsure{1|L}$-multiplicities corresponding to a given ``type'' $\Pi_0$ form a \uline{2-category} \inserted{(in groupoids)} fibered over $(\mathrm{Sch})$, and not just a fibered \uline{category}\ldots]

\scanpage{197}
\mspage{746}

We shall therefore have an exact sequence
\diagram{\begin{tikzcd}[ampersand replacement=\&, column sep=small]
(13)\quad 1 \arrow[r] \&
\widehat{\Pi}_{g,\nu} \arrow[r] \&
\underbrace{\Pi_1({\unsure{U|M}_{g,\nu}}_{\mathbb{Q}})}_{\mathcal{M}_{g,\nu}} \arrow[r] \&
\underbrace{\Pi_1({M_{g,\nu}}_{\mathbb{Q}})}_{\mathcal{N}_{g,\nu}} \arrow[r] \& 1,
\end{tikzcd}}
where ${U_{g,\nu}}_{\mathbb{Q}}$ is the ``universal curve \unsure{of|of the} type $g,\nu$'' (in characteristic $0$). Here it is $\Pi_1(U_{g,\nu}\unsure{\deleted{\mathbb{Q}}|\xi}_{\mathbb{Q}})$ which appears as a profinite $TL$-group.

I would like to bring out a property common to this ``universal'' case (13), \deleted{\unsure{and}} \deleted{(which \unsure{is} \unsure{reduced} moreover \unsure{essential-|essentially}} \inserted{and to the case \unsure{$1$-multiplicity|of $1$ multiplicity} \illegible{1} relative} to a base field $k$ an extension \uline{of finite type} of $\mathbb{Q}$ (the first case reducing moreover essentially to the second, by taking a ``generic point'' of $M_{g,\nu}$\ldots). For this, let us consider an $LT$-group $\unsure{\mathcal{M}|\mathcal{A}}$, and a \deleted{of} homomorphism
\diagram{\begin{tikzcd}[ampersand replacement=\&]
k_i:T \arrow[r] \& \deleted{\unsure{\mathcal{G}|\Pi}}\,\mathcal{M}
\end{tikzcd}}
belonging to one of the $k_i^{\unsure{\natural|\sharp}}$. One can

\scanpage{198}

for example choose a normal subgroup of definition $\Pi$ in $\mathcal{M}$, which will therefore be a group with loops, \deleted{\unsure{and one will have|and one has}} an exact sequence
\diagram{\begin{tikzcd}[ampersand replacement=\&, column sep=small]
(14)\quad 1 \arrow[r] \& \Pi \arrow[r] \& \mathcal{M} \arrow[r] \& \mathcal{N} \arrow[r] \& 1,
\end{tikzcd}}
and a homomorphism
\diagram{(15)\qquad \begin{tikzcd}[ampersand replacement=\&]
k:T \arrow[r,hook] \& \Pi
\end{tikzcd},\qquad
$k(T)\overset{\mathrm{dfn}}{=}L$
}
relative to the loop structure of $\Pi$. Let us consider the normalizer of $\unsure{L|T}$ in $\mathcal{M}$
\[
(16)\qquad \operatorname{Norm}_{\mathcal{M}}(L)\overset{\mathrm{dfn}}{=}N(L)
\]
and consider the action of $N(L)/L\subset\mathcal{N}$ \deleted{\unsure{in}} on $L$. (\uline{N.B.} One readily sees that the image of $N(L)/L$ in $\mathcal{N}$ is an open subgroup of $\mathcal{N}$\ldots) This action is by a ``cyclotomic \unsure{character}'', induced by a \unsure{character}
\diagram{(17)\qquad \begin{tikzcd}[ampersand replacement=\&]
\mathcal{N} \arrow[r,"\chi_{\mathcal{N}}"] \& \widehat{\mathbb{Z}}^{*}
\end{tikzcd}.}
The essential hypothesis is that the image

\scanpage{199}
\mspage{747}

of $\deleted{\chi}\;\mathcal{N}$ under $\chi$ (or also of $N(L)/L$ under $\chi$) is an open subgroup of $\widehat{\mathbb{Z}}^{*}$. This hypothesis has a character intrinsic to the $LT$-group $\mathcal{M}$ itself (independently of the choice of $\Pi$) --- for the cyclotomic character
\diagram{(18)\qquad \begin{tikzcd}[ampersand replacement=\&]
\chi_{\mathcal{M}}:\mathcal{M} \arrow[r] \& \widehat{\mathbb{Z}}^{*}
\end{tikzcd}}
has a character intrinsic to $\mathcal{M}$, and the hypothesis means that the image of the latter in $\widehat{\mathbb{Z}}^{*}$ is open.

Returning to $L$ and $N(L)$, one sees that the hypothesis \unsure{made|made} implies this:
\[
\begin{array}{ll}
(19)&\text{Every open subgroup of }N(L)/L\\
&\text{operates nontrivially on }L.
\end{array}
\]

One knows, by \unsure{standard|standard} arguments that this implies that for every \deleted{good} field $K$ \deleted{of characteristic $\unsure{0}$} (in particular an $\ell$-adic field $\mathbb{Q}_{\ell}$),

\scanpage{200}

and every homomorphism from $\mathcal{M}$ into a group $G(K)$; where $G$ is an affine group scheme of finite type over $K$; \deleted{\unsure{the|u}}
\diagram{(20)\qquad \begin{tikzcd}[ampersand replacement=\&]
u:\mathcal{M} \arrow[r] \& G(K)
\end{tikzcd},}
the restriction of $u$ to $L$ is ``quasi-unipotent'', i.e.
\[
\begin{array}{ll}
(21)&\exists\ \text{open subgroup }L'\text{ of }L,\text{ such that}\\
&u(L')\subset G(K)\text{ is made up of unipotent}\\
&\text{elements},
\end{array}
\]
--- which amounts to the same thing, if $\ell_i$ is a generator of $L=L_i$, that one has
\[
\begin{array}{ll}
(22)&u(\ell_i)\text{ is a quasi-}\\
&\text{unipotent element of }G(K),\text{ i.e. }\exists\ n\in\mathbb{N}^{*},\\
&\text{such that }\ell_i^{n}\text{ is unipotent}.
\end{array}
\]

\scanpage{201}
\mspage{748}

\textcircled{\(3^\circ\)}\quad \uline{Schematic envelope of a} \(T\)\unsure{-}\(L\)-\uline{profinite group} : \inserted{\unsure{\(a)\)}} \uline{The admissible homomorphisms.}

Let \(\mathcal{M}\) be a profinite \(T\)-\(L\)-group, whose structural character (18)
\diagram{\begin{tikzcd}[ampersand replacement=\&]
\mathcal{M} \arrow[r, "\chi_{\mathcal{M}}"] \& \widehat{\mathbb{Z}}^{*}
\end{tikzcd}}
has an open image, \(\Pi\) a \unsure{normal} subgroup of definition (it is therefore a profinite \deleted{\unsure{\(L\)-}}group \inserted{with loops}\unsure{)|]}, and \(\Pi_0\subset\Pi\) a \uline{discrete} subgroup of \(\Pi\), which is a ``discretification'' of it --- it is therefore a discrete \unsure{group} with loops, \unsure{etc.|and} \(\Pi\) is identified with its profinite completion\ldots

Let \(G_{\mathbb{Q}}\) be a group scheme \inserted{affine of finite type} over \(\mathbb{Q}\). Let \(\mathcal{A}_{\mathbb{Q}}\) be the ring of adeles of \(\mathbb{Q}\)
\[
(23)\qquad \mathcal{A}_{\mathbb{Q}}=\widehat{\mathbb{Z}}\otimes_{\mathbb{Z}}\mathbb{Q}\subset\prod_{\ell\in\mathbb{P}}\mathbb{Q}_{\ell}
\]
and let us consider \unsure{\deleted{the}|the} group of adelic points
\[
(24)\qquad G(\mathcal{A}_{\mathbb{Q}})\subset\prod_{\ell\in\mathbb{P}}G(\mathbb{Q}_{\ell}).
\]
If \(G_{\mathbb{Q}}\) comes from a group \unsure{scheme|scheme}

\scanpage{202}

affine of finite type \(G\) over \(\mathbb{Z}\), then we have
\[
\begin{gathered}
(25)\qquad G(\mathcal{A}_{\mathbb{Q}})\simeq
\unsure{(|\{}(x_{\ell})_{\ell\in\mathbb{P}}\in\prod_{\ell}G(\mathbb{Q}_{\ell})
\mid x_{\ell}\in G(\mathbb{Z}_{\ell})\\
\hfill\text{for almost }\unsure{\mathrm{tt}|\forall}\ell\}
\end{gathered}
\]

\deleted{One says that \unsure{a|one}}

We give ourselves in addition a homomorphism
\diagram{\begin{tikzcd}[ampersand replacement=\&]
(26)\quad G_{\mathbb{Q}} \arrow[r, "\chi_G"] \& \unsure{\mathbb{G}_{m\mathbb{Q}}|G_{m\mathbb{Q}}}.
\end{tikzcd}}
We are interested in \deleted{\unsure{a|these}}\inserted{some} group homomorphisms
\diagram{\begin{tikzcd}[ampersand replacement=\&]
(27)\quad \mathcal{M} \arrow[r] \& G(\mathcal{A}_{\mathbb{Q}})
\end{tikzcd}}
satisfying properties which we are going to enumerate.

\(1^\circ)\) \uline{Continuity}. This is equivalent to saying that \(a)\) the homomorphisms \(\mathcal{M}\longrightarrow G(\unsure{\mathbb{Q}_{\ell}})\) are \unsure{continuous|continuous}, and \(b)\) \(\exists S\) a finite subset of \(\mathbb{P}\), such that for \(\ell\notin S\), one has \deleted{\illegible{1}} a factorization (continuous)
\diagram{\begin{tikzcd}[ampersand replacement=\&]
(28)\quad \mathcal{M} \arrow[r] \& G(\widehat{\mathbb{Z}}_{\ell}) \& (\ell\in\mathbb{P}\setminus S)
\end{tikzcd}}
--- condition which indeed does not depend on the chosen ``model over \(\mathbb{Z}\)'' \(G\) of \(G_{\mathbb{Q}}\).

\uline{N.B.} (28) is a homomorphism of profinite groups.

\scanpage{203}
\mspage{749}

\uline{ReN.B.} Unless I am mistaken, property \(1^\circ)\) means that one can find an affine group scheme \(G\) of finite type over \(\mathbb{Z}\), which \unsure{defines|defines} \(G_{\mathbb{Q}}\), and such that (27) factors through \(G(\widehat{\mathbb{Z}})\).
\diagram{\begin{tikzcd}[ampersand replacement=\&]
(29)\quad \deleted{\unsure{G}}\quad \mathcal{M} \arrow[r] \& G(\widehat{\mathbb{Z}}).
\end{tikzcd}}

\(2^\circ)\) \uline{Compatibility with} \(\chi\) :

\begin{samepage}
\(\unsure{(30)}\)
\diagram{\begin{tikzcd}[ampersand replacement=\&, row sep=large, column sep=large]
\& \mathcal{M} \arrow[r] \arrow[dl, "\chi_{\mathcal{M}}"'] \& G(\mathcal{A}_{\mathbb{Q}}) \arrow[dl, "\chi_G"] \\
\widehat{\mathbb{Z}} \arrow[r, hook] \& \unsure{\mathbb{G}_m|G_m}(\mathcal{A}_{\mathbb{Q}})\simeq\mathcal{A}_{\mathbb{Q}}^{*}\subset\displaystyle\prod_{\ell\in\mathbb{P}}\mathbb{Q}_{\ell}^{*} \&
\end{tikzcd}}
\end{samepage}
is commutative.

\(3^\circ)\). \uline{Compatibility with the discretification} \(\Pi_0\).

\begin{samepage}
\((31)\)
\diagram{\begin{tikzcd}[ampersand replacement=\&, row sep=large, column sep=large]
\Pi'_0 \arrow[r, hook] \arrow[d, "\text{factorization}"] \& \Pi \arrow[r, hook] \& \mathcal{M} \arrow[d] \\
G(\mathbb{Q}) \arrow[rr, hook] \&\& G(\mathcal{A}_{\mathbb{Q}})
\end{tikzcd}}
\end{samepage},
for a suitable subgroup of finite index \(\Pi'_0\) of \(\Pi_0\).

\uline{N.B.} Changing notation if need be \unsure{we may assume|we may assume} \(\Pi'_0=\Pi_0\).

\(4^\circ)\) \uline{Minimality} (\unsure{irredundancy|irredundancy}) :

If \(G'_{\mathbb{Q}}\) is a subgroup scheme of

\scanpage{204}

\(G_{\mathbb{Q}}\), such that \(\deleted{\unsure{G}}\ \mathcal{M}\longrightarrow G(\mathcal{A}_{\mathbb{Q}})\) factors through \(G'(\mathcal{A}_{\mathbb{Q}})\), then \(G=G'\).

\uline{N.B.} If this hypothesis is not satisfied, then among all the \(G'_{\mathbb{Q}}\subset G_{\mathbb{Q}}\) which violate it, there is a \inserted{(unique)} smallest one, and on replacing \(\chi_G\) by \(\chi_{G'}=\chi_G\mid G'_{\mathbb{Q}}\ldots\), the hypotheses \(1^\circ)\) to \(4^\circ)\) are satisfied for \(G'_{\mathbb{Q}},\chi'\), and \(\mathcal{M}\longrightarrow G'(\mathcal{A}_{\mathbb{Q}})\).

\uline{Remark}. Let \(\mathcal{G}_0\subset\mathcal{M}\) be the subgroup of \(\mathcal{M}\) formed of the \(g\in\mathcal{M}\) which are such that \(\unsure{g(\Pi_0)|g\ell(\Pi_0)}\) and \(\Pi_0\) are commensurable. \unsure{That is|That is} \inserted{the \(g\)} which normalize the ``discretification germ'' \(\Pi_0\) of \(\Pi\).

\margin{\(\mathcal{G}_0/\unsure{\Pi_0}\) plays the role of a \uline{discrete} Teichmüller group}

It is possible that condition \(3^\circ\) will have to be strengthened, by requiring that we have
\diagram{\begin{tikzcd}[ampersand replacement=\&]
\mathcal{G}_0 \arrow[r] \& G(\mathbb{Q}).
\end{tikzcd}}
The future will decide it!

\scanpage{205}
\mspage{750}

Reflection \unsure{done|done}, I prefer to put things in a slightly \unsure{different|different} way, by introducing, in the \inserted{profinite} \(LT\)-group \(\mathcal{M}\), a normal \inserted{subgroup} \deleted{defining the \(LT\)-\unsure{structure|structure}} \(\Pi\) with an \(L\)-structure (not \unsure{necessari\textsuperscript{ly}|necessarily} a loop structure --- for example one can have \(\Pi=\mathrm{SL}\unsure{(}2,\mathbb{Z})\)) defining the \(LT\)-structure of \(\mathcal{M}\), a discretification \(\Pi_0\) of \(\Pi\) (in an obvious sense, for \(L\)-structures) a discrete subgroup \(\mathcal{G}_0\) \inserted{\(\supset\Pi_0\)} normalizing \(\Pi_0\) (N.B. in \unsure{important} special cases \unsure{will be} \(\mathcal{G}_0=\Pi_0\)), such that \inserted{its closure} \(\mathcal{G}_0^{\deleted{\unsure{\wedge|*}}}\) \inserted{\(\overset{\mathrm{dfn}}{=}\mathcal{G}\)} in \(\mathcal{M}\) is normal in \(\mathcal{M}\).

We thus have a diagram of groups

\begin{samepage}
\((32)\)
\diagram{\begin{tikzcd}[ampersand replacement=\&,
row sep=large, column sep=large,
execute at end picture={
  \path (s205p.north) -- coordinate[midway] (s205pgmid) (s205g.north);
  \path (s205g.north) -- coordinate[midway] (s205gmmid) (s205m.north);
  \path (s205p.north) -- coordinate[midway] (s205pmmid) (s205m.north);
  \path (s205pzero.south) -- coordinate[midway] (s205zeromid) (s205gzero.south);
  \draw ([yshift=5mm]s205p.north) .. controls ([yshift=8mm]s205p.north) and ([xshift=-2mm,yshift=5mm]s205pgmid) .. ([yshift=8mm]s205pgmid) .. controls ([xshift=2mm,yshift=5mm]s205pgmid) and ([yshift=8mm]s205g.north) .. ([yshift=5mm]s205g.north);
  \node[above] at ([yshift=8mm]s205pgmid) {$\mathcal{T}$};
  \draw ([yshift=5mm]s205g.north) .. controls ([yshift=8mm]s205g.north) and ([xshift=-2mm,yshift=5mm]s205gmmid) .. ([yshift=8mm]s205gmmid) .. controls ([xshift=2mm,yshift=5mm]s205gmmid) and ([yshift=8mm]s205m.north) .. ([yshift=5mm]s205m.north);
  \node[above] at ([yshift=8mm]s205gmmid) {$\Gamma$};
  \draw ([yshift=17mm]s205p.north) .. controls ([yshift=20mm]s205p.north) and ([xshift=-2mm,yshift=17mm]s205pmmid) .. ([yshift=20mm]s205pmmid) .. controls ([xshift=2mm,yshift=17mm]s205pmmid) and ([yshift=20mm]s205m.north) .. ([yshift=17mm]s205m.north);
  \node[above] at ([yshift=20mm]s205pmmid) {$\mathcal{N}$};
  \draw ([yshift=-3mm]s205pzero.south) .. controls ([yshift=-6mm]s205pzero.south) and ([xshift=-2mm,yshift=-3mm]s205zeromid) .. ([yshift=-6mm]s205zeromid) .. controls ([xshift=2mm,yshift=-3mm]s205zeromid) and ([yshift=-6mm]s205gzero.south) .. ([yshift=-3mm]s205gzero.south);
  \node[below] at ([yshift=-6mm]s205zeromid) {$\mathcal{T}_0$};
}]
\unsure{1} \arrow[r, hook] \& |[alias=s205p]| \underset{=\widehat{\Pi}_0}{\Pi} \arrow[r, hook] \& |[alias=s205g]| \underset{=\widehat{\mathcal{G}}_0}{\mathcal{G}^{\deleted{\unsure{\wedge|*}}}_{\deleted{\unsure{0}}}} \arrow[r, hook] \& |[alias=s205m]| \mathcal{M} \\
1 \arrow[r, hook] \& |[alias=s205pzero]| \Pi_0 \arrow[u, hook] \arrow[r, hook] \& |[alias=s205gzero]| \mathcal{G}_0 \arrow[u, hook] \&
\end{tikzcd}}
\end{samepage},

where the first row is formed of profinite groups, the second of discrete groups. We will note that \(\mathcal{G}\) is identified with a quotient of the profinite completion of \(\mathcal{G}_0\), but

\scanpage{206}

I do not require for the moment that it be isomorphic to this completion itself. \unsure{Thus} \unsure{idem|item} for $\mathcal{T}$ and $\mathcal{T}_0$. The group $\mathcal{T}_0\inserted{=\mathcal{G}_0/\Pi_0}$ plays the role of a discrete \inserted{``direct''} Teichmüller group, the group $\mathcal{T}$ that of a completion of a Teichmüller group, the group $\Gamma=\mathcal{M}/\mathcal{G}$ plays the role of a profinite Galois group. \deleted{The \unsure{situation}} I am interested mainly in two extreme situations

a) $\mathcal{G}_0=\Pi_0$ whence $\mathcal{G}=\Pi$ --- situation coming from a smooth relative \unsure{$1$-multiplicity} over a field \deleted{$k$} extension of finite type of $\mathbb{Q}$, or more generally \deleted{from a scheme}\inserted{over a} scheme \inserted{(or a multiplicity)} \unsure{integral} of finite type over $\mathbb{Q}$, \deleted{\unsure{$\Gamma$|$k$} \unsure{in fact}}\inserted{$\mathcal{N}\unsure{\xrightarrow{\sim}|\simeq}\Gamma$ is} \deleted{\unsure{is}} the fundamental group of the said \deleted{\unsure{this}} scheme or of the said multiplicity,

b) $\mathcal{G}_0/\Pi_0\simeq$ Teichmüller group of $\Pi_0$ \inserted{or of finite index in this Teichmüller group} --- in this case, one is interested mainly in a ``modular'' situation, where

\scanpage{207}
\mspage{751}

$\mathcal{N}$ is interpreted as the fundamental group of \unsure{their|the} appropriate modular multiplicity, $\mathcal{T}$ as its ``geometric'' part, $\mathcal{T}_0$ as its ``transcendental'' part, and $\Gamma$ as $\Gamma_{\unsure{\overline{\mathbb{Q}}_0|\overline{\mathbb{Q}}}/\mathbb{Q}}$ ....

We make on (32) the hypothesis that the cyclotomic character on $\mathcal{M}$ (operating on $\Pi$ by automorphisms of $L$-structures) is trivial on $\mathcal{G}$, a fortiori on $\Pi$ (thus
\[
(33)\qquad \Pi=\Pi^{+},\quad \mathcal{G}=\mathcal{G}^{+}\qquad ),
\]
i.e. that it factors as
\diagram{(34)\qquad
\begin{tikzcd}[ampersand replacement=\&, column sep=large]
\Gamma \arrow[r, "\chi_{\Gamma}"] \& \widehat{\mathbb{Z}}^{*}
\end{tikzcd},
}
\uline{whose image is an open subgroup.}

If $(G_{\mathbb{Q}},\chi_G)$ is as above (page 748), we are interested in the homomorphisms (27)
\diagram{
\begin{tikzcd}[ampersand replacement=\&, column sep=large]
\mathcal{M} \arrow[r] \& G(\mathcal{A}_{\mathbb{Q}})
\end{tikzcd}
}
satisfying the conditions $1^{\circ})$ $2^{\circ})$ $\deleted{\unsure{2}}\inserted{3}^{\circ})$ $4^{\circ})$ as above, but where $3^{\circ})$ is \unsure{formulated} in the slightly stronger form

($3^{\circ}$ strong) (Compatibility with Teichmüller discretification $\mathcal{G}_0$) : We have a \unsure{factoriza-}

\scanpage{208}

tion $\mathcal{G}_0\longrightarrow G(\mathbb{Q})$, making \unsure{thus} the diagram commutative
\diagram{(35)\qquad
\begin{tikzcd}[ampersand replacement=\&, column sep=large, row sep=large]
\mathcal{G}_0 \arrow[r, hook] \arrow[d] \& \mathcal{M} \arrow[d] \\
G(\mathbb{Q}) \arrow[r, hook] \& G(\mathcal{A}_{\mathbb{Q}})
\end{tikzcd}.
}

Structures of the form (32) probably deserve a name, I am going to call \unsure{them} --- if only provisionally, \uline{Galois-Teichmüller \unsure{structures}.}
\margin{Leaving it open later to introduce more draconian conditions on \unsure{this structure}}

If $G$ is an affine algebraic group over $\mathbb{Q}$, a homomorphism $\mathcal{M}\unsure{\longrightarrow|{-}}G(\mathcal{A}_{\mathbb{Q}})$ satisfying the conditions $1^{\circ})$ $2^{\circ})$ $3^{\circ})$ \inserted{(strong)} $4^{\circ})$ is called a ``Lie homomorphism'' of the Galois-Teichmüller structure.

I will henceforth denote by $\underline{\mathcal{M}}_{\mathbb{Q}}$ \inserted{(or \unsure{simply} $\underline{\mathcal{M}}$ ?)} the group scheme \inserted{$G_{\mathbb{Q}}$} over $\mathbb{Q}$. \deleted{\unsure{defined}}

I denote by \deleted{$\unsure{\mathcal{G}|\mathcal{M}}$} $\underline{\mathcal{G}}_{\mathbb{Q}}$ the subgroup scheme of $\underline{\mathcal{M}}_{\mathbb{Q}}$ generated by the

\scanpage{209}
\mspage{752}

image of $\mathcal{G}_0\longrightarrow\deleted{\unsure{\mathcal{M}|G}}G(\mathbb{Q})$, by $\underline{\Pi}_{\mathbb{Q}}$ the one generated by the image of $\Pi\longrightarrow G(\mathbb{Q})$

We thus have \deleted{\unsure{a}} subgroup schemes
\[
(36)\qquad \underline{\Pi}_{\mathbb{Q}}\subset
\deleted{\underline{\unsure{\mathcal{M}}}}\inserted{\underline{\mathcal{G}}}_{\mathbb{Q}}
\subset\underline{\mathcal{M}}_{\mathbb{Q}}
\]
\uline{normal} in $\underline{\mathcal{M}}_{\mathbb{Q}}$, giving rise to quotient group schemes
\[
(37)\qquad
\left\{
\begin{aligned}
&\underline{\mathcal{T}}_{\mathbb{Q}}\simeq
\underline{\mathcal{G}}_{\mathbb{Q}}/\underline{\Pi}_{\mathbb{Q}},\quad
\underline{\Gamma}_{\mathbb{Q}}\simeq
\underline{\mathcal{M}}_{\mathbb{Q}}/\underline{\mathcal{G}}_{\mathbb{Q}},\\
&\qquad\underline{\mathcal{N}}_{\mathbb{Q}}\simeq
\underline{\mathcal{M}}_{\mathbb{Q}}/\underline{\Pi}_{\mathbb{Q}}
\end{aligned}
\right.
\]
--- giving rise to four exact sequences of group schemes, the last of which is
\diagram{(38)\qquad
\begin{tikzcd}[ampersand replacement=\&, column sep=large]
1 \arrow[r] \& \underline{\mathcal{T}}_{\mathbb{Q}} \arrow[r] \&
\underline{\mathcal{N}}_{\mathbb{Q}} \arrow[r] \& \underline{\Gamma}_{\mathbb{Q}} \arrow[r] \& 1
\end{tikzcd}.
}

\deleted{I \unsure{am going to|see} \unsure{reflect|recall}}\inserted{Let us note \unsure{that}} the character $\chi_G:G_{\mathbb{Q}}\longrightarrow \unsure{\mathbb{G}_m|G_m}$ is an epimorphism, and factors through $\underline{\Gamma}_{\mathbb{Q}}$ as
\diagram{(39)\qquad
\begin{tikzcd}[ampersand replacement=\&, column sep=large]
\underline{\Gamma}_{\mathbb{Q}} \arrow[r, "\chi_{\underline{\Gamma}}"] \& \unsure{\mathbb{G}_{m\mathbb{Q}}|G_{m\mathbb{Q}}}
\end{tikzcd}.
}
